\documentclass[11pt,letterpaper]{article}
\usepackage[T1]{fontenc}
\usepackage[utf8]{inputenc}
\usepackage{lmodern}
\usepackage{amsmath,amssymb,amsthm,mathtools}
\usepackage[margin=1in,headheight=15pt]{geometry}
\usepackage{microtype,booktabs,array,longtable}
\usepackage[dvipsnames]{xcolor}
\usepackage{enumitem,fancyhdr,titlesec}
\usepackage{listings}
\usepackage[colorlinks=true,linkcolor=MidnightBlue,citecolor=MidnightBlue,urlcolor=MidnightBlue]{hyperref}
\usepackage{bookmark}
\definecolor{ink}{HTML}{17364D}
\definecolor{accent}{HTML}{1C6675}
\definecolor{pale}{HTML}{F1F5F7}
\titleformat{\section}{\large\bfseries\color{ink}}{\thesection}{0.75em}{}
\titleformat{\subsection}{\normalsize\bfseries\color{accent}}{\thesubsection}{0.65em}{}
\pagestyle{fancy}
\fancyhf{}
\fancyhead[L]{\small\color{ink}Delayed stabilization of power towers}
\fancyhead[R]{\small Research note}
\fancyfoot[C]{\thepage}
\renewcommand{\headrulewidth}{0.3pt}
\setlength{\parskip}{4pt}
\setlength{\parindent}{1.2em}
\setlength{\emergencystretch}{2em}
\linespread{1.04}
\setcounter{tocdepth}{1}
\numberwithin{equation}{section}
\newtheorem{theorem}{Theorem}[section]
\newtheorem{lemma}[theorem]{Lemma}
\newtheorem{proposition}[theorem]{Proposition}
\newtheorem{corollary}[theorem]{Corollary}
\theoremstyle{definition}
\newtheorem{definition}[theorem]{Definition}
\newtheorem{example}[theorem]{Example}
\newtheorem{question}[theorem]{Question}
\theoremstyle{remark}
\newtheorem{remark}[theorem]{Remark}
\newcommand{\N}{\mathbb N}
\newcommand{\Z}{\mathbb Z}
\newcommand{\vp}[2]{v_{#1}\!\left(#2\right)}
\newcommand{\res}[2]{\left[#1\right]_{#2}}
\newcommand{\ceil}[1]{\left\lceil #1\right\rceil}
\newcommand{\floor}[1]{\left\lfloor #1\right\rfloor}
\DeclareMathOperator{\gcdop}{gcd}
\DeclareMathOperator{\lcm}{lcm}
\lstset{basicstyle=\small\ttfamily,breaklines=true,columns=fullflexible,
 keepspaces=true,showstringspaces=false,frame=single,rulecolor=\color{accent},
 backgroundcolor=\color{pale},xleftmargin=0.4em,xrightmargin=0.4em,
 aboveskip=10pt,belowskip=10pt}
\hypersetup{pdftitle={Delayed stabilization of power towers},
 pdfsubject={A counterexample to a published congruence-speed conjecture, exact onset laws, and radix congruences},
 pdfauthor={Research report prepared with AI assistance}}

\begin{document}
\begin{titlepage}
\thispagestyle{empty}
\vspace*{0.6in}
\noindent{\color{accent}\large TETRATION \; / \; EXACT ARITHMETIC \; / \; CONJECTURE TESTING}\par
\vspace{0.35in}
\noindent{\color{ink}\Huge\bfseries Delayed stabilization\\[5pt] of power towers}\par
\vspace{0.25in}
\noindent{\Large A counterexample to a published early-speed conjecture,\\[4pt]
exact onset laws, and sharp radix congruences}\par
\vspace{0.35in}
\noindent September 19, 2026 \hfill Version 1.0\par
\vspace{0.35in}
\begin{abstract}
We investigate the height at which the number of newly stabilized trailing digits of an integer power tower becomes constant. A conjecture of Rip\`a predicts that, for every prime base ending in 9, this happens by height 2. With the height convention of the source, we prove that the smallest prime counterexample is $2749$: its congruence speeds are $3,3,2,2,\ldots$. More generally, if $a\equiv9\pmod{10}$, $s=v_2(a-1)$, $t=v_2(a+1)$, $u=s+t-1$, and $e=v_5(a+1)$, then the number of stable decimal digits at height $h$ is exactly $\min(s+hu,he)$. This determines the entire transient and its exact endpoint. A Chinese-remainder construction and Dirichlet's theorem give infinitely many prime bases with each prescribed onset height $H\ge3$. Among integers ending in 9, we also obtain the exact tail density $4/[9(5\cdot10^{K-1}-1)]$ for $H>K$, $K\ge2$.

A complementary theorem proves that, for even $B\ge4$ and tower base $B-1$, consecutive height-$h$ towers differ by $-2B^h$ modulo $B^{h+1}$. It resolves three conjectural comments in OEIS A324017: two are true, while the third fails in every nondegenerate case. We give complete proofs, an exact digit-lifting algorithm, a restricted general-radix extension, and a consequence for higher up-arrow operations. The accompanying dependency-free Python package passes 118,271 exact assertions.
\end{abstract}
\vfill
\noindent\colorbox{pale}{\parbox{0.95\textwidth}{\small
\textbf{Status and scope.} This is an unrefereed research note, not a claim of verified priority. The eventual-speed formula is already in the literature; the focus here is the exact transient, the explicit conjecture refutation, and its consequences. The conjecture remains labeled as such in the source version examined, but the literature search does not certify that no earlier resolution exists. All asserted results are proved below; computational tests are supplementary, not proofs of infinite statements.}}
\end{titlepage}

\tableofcontents
\newpage

\section{The selected problem and its provenance}\label{sec:problem}

Integer tetration provides a useful contrast between size and information. The ordinary value of a tower rapidly becomes impossible to write down, but its residue modulo a fixed integer can settle after only a few levels. The present problem concerns not merely whether trailing digits stabilize, but \emph{when the rate of their stabilization becomes permanent}.

We choose a precise published conjecture rather than pose an unverified question of our own. Rip\`a's \emph{The congruence speed formula} \cite{Ripa2021} introduces the prime set
\begin{equation}\label{eq:source-set}
 \widehat{\mathbb A}=\{(k+1)10^n-1:\ n\ge1,\ k\ge0,\ (k+1)10^n-1\text{ is prime}\}.
\end{equation}
This is exactly the set of primes congruent to $9$ modulo $10$: the forward implication is immediate, and the reverse implication uses $n=1$ and $k=(p+1)/10-1$.

The conjecture is numbered \textbf{Conjecture 2 on journal page 57}, and \textbf{Conjecture 4.1 in arXiv:2208.02622v1, page 14}. In the notation defined carefully in Section~\ref{sec:notation}, it says
\begin{equation}\label{eq:conjecture}
  p\in\widehat{\mathbb A},\ h\ge2
  \quad\Longrightarrow\quad V(p,h)=V(p).
\end{equation}
Its point is a \emph{uniform onset bound}, not an assertion that the limiting speed exists. Confusing these statements would miss the counterexample.

\subsection{What was already known, and what is established here}

Rip\`a and Onnis \cite{RipaOnnis2022}, especially Corollary 2.2 and Section 2.2, already give the eventual decimal speed in terms of 2-adic and 5-adic valuations and discuss bounds on transient behavior. Their displayed example $a=781249$ has a nontrivial transient; it is composite, since $781249=7\cdot233\cdot479$, so that example alone does not refute the prime conjecture. The present note does \emph{not} claim to discover the eventual-speed formula or the general phenomenon of delayed stabilization.

The results proved here are an explicit minimal prime counterexample, a necessary-and-sufficient condition for every base ending in 9, a closed formula for the exact onset, prime families with unbounded onset, and a density formula for the delays. A second investigation addresses the three conjectural comments in the retrieved version of OEIS A324017 \cite{OEIS324017}.

The primary-source audit was performed on September 19, 2026. It included the journal and arXiv versions of \cite{Ripa2021}, the relevant sections of \cite{RipaOnnis2022}, and the current A324017 entry. Targeted searches for the conjecture and for the base $2749$ did not locate an explicit prior resolution. That is a search outcome, not a guarantee of current open status or worldwide novelty. Mathematical correctness of the refutation below does not depend on priority.

\subsection{An overview of the conclusions}

For $a\equiv9\pmod{10}$, put $s=v_2(a-1)$, $t=v_2(a+1)$, $u=s+t-1$, and $e=v_5(a+1)$. The essential formulas are
\begin{align}
 C_a(h)&=\min\{s+hu,he\},\label{eq:overview-count}\\
 V(a)&=\min\{u,e\},\label{eq:overview-speed}\\
 H(a)&=\begin{cases}
 1,&e\le u,\\
 1+\ceil{s/(e-u)},&e>u.
 \end{cases}\label{eq:overview-onset}
\end{align}
Here $C_a(h)$ counts stable decimal digits at height $h$, and $H(a)$ is the first height from which the speed always equals its eventual value. The proof is an elementary comparison of two exact affine valuation laws. Primality plays no role until the arithmetic-progression argument.

In particular,
\[
 H(a)>2\quad\Longleftrightarrow\quad
 a\equiv1\pmod4\ \text{ and }\ s<e<2s.
\]
The complete answer therefore distinguishes the two classes $9$ and $19$ modulo $20$. Every base in the latter class satisfies the proposed bound; some bases in the former do not.

\section{Definitions and the height-indexing issue}\label{sec:notation}

Throughout the main argument, $a$ is an odd integer at least $3$. Define
\begin{equation}\label{eq:towers}
 T_0(a)=1,\qquad T_{h+1}(a)=a^{T_h(a)}\quad(h\ge0),
 \qquad D_h(a)=T_{h+1}(a)-T_h(a).
\end{equation}
Thus $T_h(a)=a\uparrow\uparrow h$ and $D_0(a)=a-1$. Every $T_h(a)$ is odd. The towers increase strictly, so every $D_h(a)$ is a positive even integer.

For a prime $q$ and a nonzero integer $N$, $v_q(N)$ is the largest $j\ge0$ such that $q^j\mid N$. For any modulus $M\ge1$, $[N]_M$ means the least nonnegative residue of $N$ modulo $M$.

\begin{definition}[Stable count, speed, and onset]\label{def:counts}
For $a\equiv9\pmod{10}$ and $h\ge0$, let
\begin{equation}\label{eq:Cdef}
 C_a(h)=\min\{v_2(D_h(a)),v_5(D_h(a))\}.
\end{equation}
Equivalently, $10^{C_a(h)}$ divides $D_h(a)$ but $10^{C_a(h)+1}$ does not. For $h\ge1$, define
\begin{equation}\label{eq:Vdef}
 V(a,h)=C_a(h)-C_a(h-1).
\end{equation}
When $V(a,h)$ is eventually constant, call that constant $V(a)$ and put
\begin{equation}\label{eq:Hdef}
 H(a)=\min\{H\ge1:\ V(a,h)=V(a)\text{ for every }h\ge H\}.
\end{equation}
\end{definition}

For bases ending in 9, $C_a(0)=0$ because $5\nmid a-1$. Thus there is no initial additive offset hidden in the speed convention. In the source definition, height $h$ compares the common digits of $(T_{h-1},T_h)$ with those of $(T_h,T_{h+1})$. This is precisely \eqref{eq:Vdef}, not $C_a(h+1)-C_a(h)$. The counterexample is therefore not an indexing artifact.

Initially, $C_a(h)$ is defined through \emph{two consecutive} towers. Corollary~\ref{cor:persistence} proves that these digits remain unchanged at every later height. This justifies the word ``stable'' without assuming it in advance.

An onset height is also distinct from the first height at which \emph{a specified number of digits} has stabilized. For example, $H(2749)=3$ says that the gain per level is permanent from height 3; it does not say that every desired number of digits is already fixed there.

\section{Elementary valuation laws}\label{sec:local}

For completeness, the required lifting-the-exponent identities are proved here. This avoids making the main result depend on an unexplained computational rule.

\begin{lemma}[Odd-prime lifting]\label{lem:oddLTE}
Let $q$ be an odd prime, let $q\mid A-1$, and let $N\ge1$. Then
\[
 v_q(A^N-1)=v_q(A-1)+v_q(N).
\]
\end{lemma}
\begin{proof}
First suppose $q\nmid N$. Factoring $A^N-1$ gives
\[
 \frac{A^N-1}{A-1}=1+A+\cdots+A^{N-1}\equiv N\not\equiv0\pmod q,
\]
so no additional factor of $q$ occurs.

Next write $A=1+q^r c$ with $r\ge1$ and $q\nmid c$. In the binomial expansion of $A^q-1$, the linear term is $q^{r+1}c$. Terms of index $2\le j\le q-1$ have valuation at least $rj+1\ge r+2$, and the last term has valuation $rq\ge r+2$. The unique term of smallest valuation is therefore the linear term, giving $v_q(A^q-1)=r+1$.

Write $N=q^k m$ with $q\nmid m$. Apply the first step to $A^m$, and then the second step $k$ times. The valuation increases by exactly $k$.
\end{proof}

\begin{lemma}[Binary lifting]\label{lem:twoLTE}
For odd $a$ and even $N\ge2$,
\begin{equation}\label{eq:twoLTE}
 v_2(a^N-1)=v_2(a-1)+v_2(a+1)+v_2(N)-1.
\end{equation}
\end{lemma}
\begin{proof}
Write $N=2^k m$, where $k\ge1$ and $m$ is odd. For odd $m$, the quotients
\[
 \frac{a^m-1}{a-1}=1+a+\cdots+a^{m-1},\qquad
 \frac{a^m+1}{a+1}=a^{m-1}-a^{m-2}+\cdots-a+1
\]
are odd. Therefore $v_2(a^m-1)=v_2(a-1)$ and $v_2(a^m+1)=v_2(a+1)$. Set $b=a^m$. Factor
\[
 b^{2^k}-1=(b-1)(b+1)\prod_{j=1}^{k-1}(b^{2^j}+1).
\]
The square of every odd integer is $1$ modulo $8$, so each factor in the product has valuation exactly $1$. Adding valuations gives \eqref{eq:twoLTE}.
\end{proof}

\begin{theorem}[Exact local affine laws]\label{thm:local}
Let $a\ge3$ be odd and write
\begin{equation}\label{eq:stu}
 s=v_2(a-1),\qquad t=v_2(a+1),\qquad u=s+t-1.
\end{equation}
For every $h\ge0$,
\begin{equation}\label{eq:binary-line}
 v_2(D_h(a))=s+hu.
\end{equation}
If $q$ is an odd prime dividing $a-1$, and $r=v_q(a-1)$, then
\begin{equation}\label{eq:minus-line}
 v_q(D_h(a))=(h+1)r.
\end{equation}
If $q$ is an odd prime dividing $a+1$, and $r=v_q(a+1)$, then
\begin{equation}\label{eq:plus-line}
 v_q(D_h(a))=hr.
\end{equation}
\end{theorem}
\begin{proof}
For $h\ge1$, subtract the two successive powers:
\begin{equation}\label{eq:Drec}
 D_h=a^{T_h}-a^{T_{h-1}}
     =a^{T_{h-1}}(a^{D_{h-1}}-1).
\end{equation}
All primes in the statement are coprime to $a$, so the first factor contributes no valuation. Since $D_{h-1}$ is even, Lemma~\ref{lem:twoLTE} gives
\[
 v_2(D_h)=v_2(D_{h-1})+s+t-1.
\]
The initial value is $v_2(D_0)=s$, proving \eqref{eq:binary-line}.

If $q\mid a-1$, Lemma~\ref{lem:oddLTE} gives
$v_q(D_h)=v_q(D_{h-1})+r$. Its initial value is $r$, proving \eqref{eq:minus-line}.

If $q\mid a+1$, then $q\nmid a-1$, because a common odd divisor of $a-1$ and $a+1$ would divide $2$. Since $D_{h-1}$ is even, apply Lemma~\ref{lem:oddLTE} to
\[
 a^{D_{h-1}}=(a^2)^{D_{h-1}/2}.
\]
Here $v_q(a^2-1)=r$ and $v_q(D_{h-1}/2)=v_q(D_{h-1})$. Again the increment is $r$, but now the initial value is $0$. This proves \eqref{eq:plus-line}.
\end{proof}

\begin{corollary}[Permanent, not accidental, agreement]\label{cor:persistence}
For any prime $q$ covered by Theorem~\ref{thm:local} and integers $k>h\ge0$,
\begin{equation}\label{eq:persistent}
 v_q(T_k(a)-T_h(a))=v_q(D_h(a)).
\end{equation}
In particular, for $a\equiv9\pmod{10}$, the first $C_a(h)$ trailing decimal digits of $T_h(a)$ agree with those of every later tower, and the next digit never agrees with any later tower.
\end{corollary}
\begin{proof}
The identity
\[
 T_k-T_h=D_h+D_{h+1}+\cdots+D_{k-1}
\]
has terms with strictly increasing $q$-adic valuations. In a sum with a unique term of smallest valuation, the valuation of the sum is that smallest value: divide by its power of $q$ and reduce modulo $q$. This proves \eqref{eq:persistent}. Apply it to $q=2,5$ and take the minimum.
\end{proof}

\section{The exact onset for every base ending in 9}\label{sec:onset}

\begin{theorem}[Complete decimal transient]\label{thm:onset}
Let $a\equiv9\pmod{10}$, and let $s,t,u$ be as in \eqref{eq:stu}. Set $e=v_5(a+1)\ge1$. Then
\begin{equation}\label{eq:exactC}
 \boxed{\ C_a(h)=\min\{s+hu,he\}\quad(h\ge0).\ }
\end{equation}
If $e\le u$, then $V(a,h)=e$ for every $h\ge1$ and $H(a)=1$.

If $e>u$, put $\delta=e-u$. In this case
\begin{align}
 C_a(h)&=uh+\min\{s,h\delta\},\label{eq:crossing}\\
 V(a,h)&=u+\min\{s,h\delta\}-\min\{s,(h-1)\delta\},\label{eq:exactspeed}\\
 H(a)&=1+\ceil{s/\delta}.\label{eq:exactH}
\end{align}
In all cases $V(a)=\min\{u,e\}$.
\end{theorem}
\begin{proof}
Since $5\mid a+1$, Theorem~\ref{thm:local} gives $v_5(D_h)=he$ and $v_2(D_h)=s+hu$. This proves \eqref{eq:exactC}.

If $e\le u$, then $he\le hu<s+hu$, so $C_a(h)=he$ and every speed is $e$.

Suppose $e>u$. Subtract $uh$ from both arguments of the minimum in \eqref{eq:exactC}; the result is \eqref{eq:crossing}, and taking its first difference gives \eqref{eq:exactspeed}. Put $r=\ceil{s/\delta}$. For $h\ge r+1$, both minima in \eqref{eq:exactspeed} equal $s$, so the speed is $u$. At $h=r$, the second minimum is $(r-1)\delta<s$, whereas the first is $s$. Hence $V(a,r)>u$, proving that $r+1$ is the \emph{first permanent} height. This establishes \eqref{eq:exactH} and the limiting speed.
\end{proof}

The minimum of two affine lines explains the delayed decrease. The 2-adic line begins with a positive intercept $s$, while the 5-adic line begins at $0$. Even when the binary slope $u$ is smaller, the 5-adic line can control several initial levels. The eventual speed records only the smaller slope; it discards the intercept responsible for the delay.

\begin{corollary}[Exact criterion for failure of the height-2 bound]\label{cor:criterion}
For $a\equiv9\pmod{10}$,
\begin{equation}\label{eq:failcriterion}
 H(a)>2
 \quad\Longleftrightarrow\quad
 a\equiv1\pmod4\ \text{ and }\ v_2(a-1)<v_5(a+1)<2v_2(a-1).
\end{equation}
Every $a\equiv19\pmod{20}$ satisfies $H(a)\le2$.
\end{corollary}
\begin{proof}
Consecutive even integers $a-1$ and $a+1$ have exactly one member divisible by $4$. If $a\equiv3\pmod4$, then $s=1$, $t\ge2$, and $u=t$. The nonconstant case has $H=1+\ceil{1/(e-u)}=2$. Thus failure is impossible in this class.

If $a\equiv1\pmod4$, then $s\ge2$, $t=1$, and $u=s$. Theorem~\ref{thm:onset} gives $H>2$ precisely when $e>s$ and $\ceil{s/(e-s)}>1$. The latter inequality is equivalent to $s>e-s$, or $e<2s$. This is \eqref{eq:failcriterion}.
\end{proof}

\begin{corollary}[A base-dependent upper bound]\label{cor:upperH}
For every $a\equiv9\pmod{10}$,
\[
 H(a)\le v_2(a-1)+1\le1+\log_2(a-1).
\]
The first bound is attained by infinitely many prime bases for every value $v_2(a-1)\ge2$.
\end{corollary}
\begin{proof}
In the nonconstant case, $\delta\ge1$ in \eqref{eq:exactH}; the constant case is immediate. The second bound follows from $2^{v_2(a-1)}\le a-1$. Attainment follows from Theorem~\ref{thm:prime-family} below.
\end{proof}

The eventual speed in Theorem~\ref{thm:onset} agrees with the previously known formulas in \cite{RipaOnnis2022}. The additional information is the full finite-height sequence and its exact stopping point, including the equality case when the two valuation lines meet at an integer height.

\section{The smallest prime counterexample}\label{sec:2749}

\begin{theorem}\label{thm:2749}
The smallest prime counterexample to \eqref{eq:conjecture} is $p=2749$. Its stable counts and speeds are
\[
 C_p(h):\quad 0,3,6,8,10,12,\ldots,
 \qquad
 V(p,h):\quad 3,3,2,2,2,\ldots.
\]
Thus $V(2749,2)=3\ne2=V(2749)$ and $H(2749)=3$.
\end{theorem}
\begin{proof}
The factorizations
\[
 2749-1=2^2\cdot3\cdot229,\qquad
 2749+1=2\cdot5^3\cdot11
\]
give $s=2$, $t=1$, $u=2$, and $e=3$. Consequently
\begin{equation}\label{eq:2749}
 C_{2749}(h)=\min\{2+2h,3h\}.
\end{equation}
The asserted counts, speeds, and onset follow immediately.

It remains to certify primality and minimality. Since $\sqrt{2749}<53$, trial division by the primes below $53$ is sufficient. The remainders are shown in Table~\ref{tab:primecert}; none is zero, so $2749$ is prime. Moreover $2749=275\cdot10-1$, so it belongs to \eqref{eq:source-set} with $n=1$ and $k=274$.

By Corollary~\ref{cor:criterion}, any smaller counterexample $a$ must have $s\ge2$ and $e\ge s+1\ge3$. Hence $250\mid a+1$. Write $a=250k-1$. The inequality $a<2749$ gives $1\le k\le10$, and $a\equiv1\pmod4$ forces $k$ odd. There are only five candidates, listed in Table~\ref{tab:minimality}. Only $749$ and $1749$ satisfy the failure criterion, and both are composite. No smaller prime counterexample exists.
\end{proof}

\begin{table}[htbp]
\centering
\caption{An elementary primality certificate for $2749$.}\label{tab:primecert}
\begin{tabular}{rrrrrrrrr}
\toprule
$q$&2&3&5&7&11&13&17&19\\
$2749\bmod q$&1&1&4&5&10&6&12&13\\
\midrule
$q$&23&29&31&37&41&43&47&\\
$2749\bmod q$&12&23&21&11&2&40&23&\\
\bottomrule
\end{tabular}
\end{table}

\begin{table}[htbp]
\centering
\caption{All possible smaller bases after the necessary congruence restrictions.}\label{tab:minimality}
\begin{tabular}{rrrrl}
\toprule
$a$&$s=v_2(a-1)$&$e=v_5(a+1)$&$H(a)$&Reason it is not a prime counterexample\\
\midrule
249&3&3&1&No delay beyond height 2\\
749&2&3&3&$749=7\cdot107$\\
1249&5&4&1&No delay beyond height 2\\
1749&2&3&3&$1749=3\cdot11\cdot53$\\
2249&3&3&1&No delay beyond height 2\\
\bottomrule
\end{tabular}
\end{table}

\subsection{A numerical certificate without enormous integers}

The main proof needs no explicit tower evaluation. Nevertheless, Table~\ref{tab:2749trace} provides an independent, finite arithmetic check. All residues are computed modulo $10^{16}$. For the displayed differences, both relevant valuations are below $16$, so they are determined exactly by these residues.

\begin{table}[htbp]
\centering\small
\caption{An exact modular trace for the base $2749$. Leading zeros are shown where useful.}\label{tab:2749trace}
\begin{tabular}{rllrrr}
\toprule
$h$&$[T_h]_{10^{16}}$&$[D_h]_{10^{16}}$&$v_2(D_h)$&$v_5(D_h)$&$C(h)$\\
\midrule
1&0000000000002749&3851756492182000&4&3&3\\
2&3851756492184749&6009668063000000&6&6&6\\
3&9861424555184749&9239329500000000&8&9&8\\
4&9100754055184749&8421750000000000&10&12&10\\
\bottomrule
\end{tabular}
\end{table}

The factor $10^6$ in $D_2$ does not imply that the gain at height 2 is $6$: one must subtract the previous stable count $C(1)=3$. This is another place where a careless reading of ``speed'' could obscure the counterexample.

\subsection{A second explicit delay}

For $a=16249$,
\[
 a-1=2^3\cdot3\cdot677,\qquad a+1=2\cdot5^4\cdot13.
\]
The included deterministic trial-division code certifies that $16249$ is prime. Its stable counts are $0,4,8,12,15,18,\ldots$, and its speeds are $4,4,4,3,3,\ldots$. Thus $H(16249)=4$. The next section explains why such examples cannot be controlled by any universal constant height.

\section{Infinitely many primes with every delayed onset}\label{sec:families}

\begin{theorem}[Unbounded prime onset]\label{thm:prime-family}
For every integer $s\ge2$, there are infinitely many primes $p\equiv9\pmod{10}$ such that
\begin{equation}\label{eq:prime-speed}
 V(p,h)=\begin{cases}
 s+1,&1\le h\le s,\\
 s,&h\ge s+1.
 \end{cases}
\end{equation}
Consequently, every integer $H\ge3$ occurs as the exact onset $H(p)$ for infinitely many prime bases ending in 9.
\end{theorem}
\begin{proof}
Impose the two congruences
\begin{equation}\label{eq:CRT}
 \begin{split}
 p&\equiv 1+2^s\pmod{2^{s+1}},\\
 p&\equiv 5^{s+1}-1\pmod{5^{s+2}}.
 \end{split}
\end{equation}
The moduli are coprime. By the Chinese remainder theorem, these conditions define one residue class $r_s$ modulo
\[
 M_s=2^{s+1}5^{s+2}.
\]
Every integer in the class is odd and is $-1$ modulo $5$, so $\gcd(r_s,M_s)=1$. Dirichlet's theorem on arithmetic progressions \cite{Dirichlet} therefore supplies infinitely many primes in this class.

The first congruence in \eqref{eq:CRT} gives $v_2(p-1)=s$ exactly; since $s\ge2$, it also gives $v_2(p+1)=1$. The second congruence gives $v_5(p+1)=s+1$ exactly. Hence $u=s$, $e=s+1$, and $\delta=1$ in Theorem~\ref{thm:onset}. Equation~\eqref{eq:exactspeed} reduces to \eqref{eq:prime-speed}, and $H(p)=s+1$.
\end{proof}

The infinitude assertion uses Dirichlet's theorem, not merely the existence of one CRT residue. Coprimality with the modulus is checked explicitly; omitting that check would invalidate the prime-family argument.

For a concrete construction of the residue, let $m_2=2^{s+1}$, $m_5=5^{s+2}$, $r_2=1+2^s$, and $r_5=5^{s+1}-1$. Then
\begin{equation}\label{eq:explicitCRT}
 r_s=\res{r_2+m_2\res{(r_5-r_2)m_2^{-1}}{m_5}}{m_2m_5},
\end{equation}
where the inverse is taken modulo $m_5$. Table~\ref{tab:CRT} gives small examples. The residue class chosen for $s=2$ does not contain $2749$; it is only one of the possible classes with the required exact valuations.

\begin{table}[htbp]
\centering
\caption{Prime families with $H=s+1$. The example primes are certified by trial division in the code.}\label{tab:CRT}
\begin{tabular}{rrrrr}
\toprule
$s$&$H$&$r_s$&$M_s$&A prime in the class\\
\midrule
2&3&749&5000&5749\\
3&4&16249&50000&16249\\
4&5&331249&500000&331249\\
5&6&2281249&5000000&2281249\\
6&7&20781249&50000000&1420781249\\
\bottomrule
\end{tabular}
\end{table}

\begin{remark}[What unboundedness means]
For each prescribed finite height $H$, a different arithmetic progression produces infinitely many prime examples. No single fixed base has an infinite transient. Indeed every fixed base in this class has the finite onset in \eqref{eq:exactH}. Unboundedness across bases and eventual constancy for each base are fully compatible.
\end{remark}

\section{How frequent are delayed bases?}\label{sec:density}

The construction proves existence, but it does not describe frequency. Here frequency is measured among \emph{all positive integers ending in 9}, not among primes. For a subset $E$ of that progression, define its relative natural density, when the limit exists, by
\[
 \operatorname{dens}_9(E)=
 \lim_{X\to\infty}
 \frac{\#\{a\le X:\ a\equiv9\pmod{10},\ a\in E\}}
 {\#\{a\le X:\ a\equiv9\pmod{10}\}}.
\]

\begin{theorem}[Exact tail of the onset distribution]\label{thm:density}
For every integer $K\ge2$, the relative natural density exists and is
\begin{equation}\label{eq:density}
 \boxed{\ \operatorname{dens}_9\{a:H(a)>K\}
 =\frac{4}{9(5\cdot10^{K-1}-1)}.\ }
\end{equation}
In particular, the relative density of bases violating $H\le2$ is $4/441$.
\end{theorem}
\begin{proof}
If $H(a)>K\ge2$, Corollary~\ref{cor:criterion} forces $a\equiv1\pmod4$, so $s\ge2$, $t=1$, and $u=s$. Put $d=e-s$. From \eqref{eq:exactH},
\[
 H(a)>K
 \quad\Longleftrightarrow\quad
 d\ge1\quad\text{and}\quad s>(K-1)d.
\]
For fixed $s\ge2$ and $e\ge1$, exact binary valuation $v_2(a-1)=s$ describes one residue among the $2^s$ odd classes modulo $2^{s+1}$. Its relative proportion is $2^{-s}$. Conditional on $a\equiv-1\pmod5$, exact valuation $v_5(a+1)=e$ has proportion $4/5^e$. The Chinese remainder theorem makes these two conditions independent in their finite residue-class counts. Thus the joint relative density is
\[
 \frac{1}{2^s}\frac{4}{5^e}.
\]
Substitute $e=s+d$ and sum over the allowed pairs:
\begin{align*}
 \sum_{d\ge1}\sum_{s\ge(K-1)d+1}
     \frac{4}{2^s5^{s+d}}
 &=\sum_{d\ge1}\frac{4}{5^d}
     \frac{10^{-((K-1)d+1)}}{1-10^{-1}}\\
 &=\frac49\sum_{d\ge1}
     \left(\frac{1}{5\cdot10^{K-1}}\right)^d\\
 &=\frac{4}{9(5\cdot10^{K-1}-1)}.
\end{align*}

To justify that this series is a natural density, rather than a formal probability calculation, truncate to $s<S$. Since $s>(K-1)d$, only finitely many $d$ then occur, and every retained set is a finite union of residue classes. The omitted part is contained in $\{a:v_2(a-1)\ge S\}$, whose relative density is $2^{1-S}$. Its upper relative density tends to zero as $S\to\infty$. Hence the finite density sums converge to the density of the full set, proving \eqref{eq:density}.
\end{proof}

For each $h\ge3$, subtraction of successive tails gives
\begin{equation}\label{eq:pointmass}
 \operatorname{dens}_9\{a:H(a)=h\}
 =\frac{4}{9(5\cdot10^{h-2}-1)}
  -\frac{4}{9(5\cdot10^{h-1}-1)}.
\end{equation}
The tail decays asymptotically as $\frac8{9}\,10^{-K}$. Long delays are therefore possible but sparse in this precise arithmetic sense.

\begin{table}[htbp]
\centering
\caption{Exact integer counts among the 100,000 positive integers below $10^6$ ending in 9. Finite counts are not used to prove the density theorem.}\label{tab:density}
\begin{tabular}{rrr}
\toprule
$K$&Number with $H(a)>K$&Limiting relative density\\
\midrule
2&908&$4/441$\\
3&90&$4/4491$\\
4&9&$4/44991$\\
5&1&$4/449991$\\
6&0&$4/4499991$\\
\bottomrule
\end{tabular}
\end{table}

No prime-density statement is needed or asserted here. Although each CRT family contains infinitely many primes, Dirichlet's infinitude theorem alone is not a density theorem for the union of these families.

\section{A general-radix extension with exact thresholds}\label{sec:general-radix}

The local argument is not inherently decimal. It works cleanly whenever the prime divisors of the chosen radix already divide $a^2-1$. This is a useful extension, but its restriction should not be silently dropped.

\begin{theorem}[Affine local laws in a supported radix]\label{thm:general-radix}
Let $a\ge3$ be odd, let $R\ge2$, and assume that every prime divisor of $R$ divides $a^2-1$. Write $R=\prod_{q\mid R}q^{f_q}$. For each such prime, define $(\alpha_q,\beta_q)$ by
\[
 (\alpha_q,\beta_q)=
 \begin{cases}
 (s,s+t-1),&q=2,\\
 (r,r),&q\text{ odd},\ q\mid a-1,\ r=v_q(a-1),\\
 (0,r),&q\text{ odd},\ q\mid a+1,\ r=v_q(a+1).
 \end{cases}
\]
The number of permanently stable radix-$R$ digits of $T_h(a)$ is
\begin{equation}\label{eq:radixC}
 C_R(h)=\min_{q\mid R}\floor{\frac{\alpha_q+h\beta_q}{f_q}}.
\end{equation}
For $m\ge1$, the first height, allowing height $0$, at which the residue modulo $R^m$ is permanent is exactly
\begin{equation}\label{eq:radixthreshold}
 J_R(m)=\max_{q\mid R}\max\left\{0,
       \ceil{\frac{mf_q-\alpha_q}{\beta_q}}\right\}.
\end{equation}
\end{theorem}
\begin{proof}
Theorem~\ref{thm:local} gives $v_q(D_h)=\alpha_q+h\beta_q$, with $\beta_q>0$. Divisibility by $R^m$ is equivalent to all inequalities $\alpha_q+h\beta_q\ge mf_q$. This proves \eqref{eq:radixC} and \eqref{eq:radixthreshold}. Corollary~\ref{cor:persistence} makes the agreement permanent, and shows that any smaller height fails to have the limiting residue modulo $R^m$.
\end{proof}

The symbol $v_R$ for composite $R$ is sometimes used informally for the largest $j$ such that $R^j$ divides an integer. This function is not generally a valuation: its product rule can fail. Formula~\eqref{eq:radixC} deliberately uses genuine prime valuations and an explicit minimum instead.

\begin{proposition}[Eventually periodic speed in a general radix]\label{prop:periodic}
Under the hypotheses of Theorem~\ref{thm:general-radix}, put
\[
 \rho=\min_{q\mid R}\frac{\beta_q}{f_q},\qquad
 \eta=\min_{\substack{q\mid R\\ \beta_q/f_q=\rho}}
             \frac{\alpha_q}{f_q}.
\]
For all sufficiently large $h$,
\[
 C_R(h)=\floor{\rho h+\eta}.
\]
The differences $C_R(h)-C_R(h-1)$ are eventually periodic. They are eventually constant if and only if $\rho$ is an integer.
\end{proposition}
\begin{proof}
For a finite collection of real numbers, taking the minimum commutes with the floor. Among the finitely many affine functions
$(\alpha_q+h\beta_q)/f_q$, every larger-slope line eventually lies above every minimum-slope line. Among the latter, the least intercept is $\eta$. This proves the displayed formula.

Write $\rho=P/Q$ in lowest terms. Then
$\floor{\rho(h+Q)+\eta}=\floor{\rho h+\eta}+P$, so the first differences have eventual period dividing $Q$. If $\rho$ is an integer, the difference is $\rho$. Conversely, an eventually constant integer difference would force $\lim C_R(h)/h$ to be an integer, while that limit is $\rho$.
\end{proof}

For instance, $a=3$ and $R=16$ give
\[
 C_{16}(h)=\floor{\frac{1+2h}{4}},
\]
whose speeds alternate $0,1,0,1,\ldots$. Thus constant speed is not a radix-independent phenomenon. This example does not contradict the decimal formulas, which use a different digit unit.

\section{A sharp one-digit law for towers of base \texorpdfstring{$B-1$}{B-1}}\label{sec:sharp}

Now let $B\ge4$ be even and put $a=B-1$. The following congruence is stronger than merely saying that one more radix-$B$ digit stabilizes at each level.

\begin{theorem}[Exact leading unstable digit]\label{thm:sharp}
For every integer $h\ge0$,
\begin{equation}\label{eq:sharp}
 \boxed{\ D_h(B-1)\equiv-2B^h\pmod{B^{h+1}}.\ }
\end{equation}
Consequently, $B^h\mid D_h$ but $B^{h+1}\nmid D_h$. For every $k>h$,
\begin{equation}\label{eq:sharp-telescope}
 T_k(B-1)-T_h(B-1)\equiv-2B^h\pmod{B^{h+1}}.
\end{equation}
The first permanent height modulo $B^m$ is exactly $m$.
\end{theorem}

We first isolate a binomial estimate that is valid for composite $B$.

\begin{lemma}\label{lem:binomial}
If $B\ge4$ is even, $r\ge0$, and $N$ is a positive multiple of $2B^r$, then
\begin{equation}\label{eq:binomial}
 (1-B)^N\equiv1-BN\pmod{B^{r+2}}.
\end{equation}
\end{lemma}
\begin{proof}
The quadratic binomial term is
$\binom N2 B^2=N(N-1)B^2/2$, which is divisible by $B^{r+2}$. For $j\ge3$, let $q\mid B$ be prime and put $f=v_q(B)$. The identity
\[
 \binom Nj=\frac Nj\binom{N-1}{j-1}
\]
gives
\[
 v_q\!\left(\binom Nj B^j\right)
 \ge rf+jf-v_q(j).
\]
For every $j\ge3$, $v_q(j)\le j-2$. One way to see this is that
$2^{j-1}>j$, so $q^{j-1}>j$ as well. Therefore the last valuation is at least $(r+2)f$. Every higher term is divisible by $B^{r+2}$, proving \eqref{eq:binomial}.
\end{proof}

\begin{proof}[Proof of Theorem~\ref{thm:sharp}]
At height $0$, $D_0=B-2\equiv-2\pmod B$, and $D_0$ is even. Suppose $h\ge1$ and the assertion holds for $h-1$. Since $B$ is even, it implies that $D_{h-1}$ is a multiple of $2B^{h-1}$. Because this exponent is even,
\[
 a^{D_{h-1}}=(B-1)^{D_{h-1}}=(1-B)^{D_{h-1}}.
\]
Lemma~\ref{lem:binomial}, with $r=h-1$, shows
\[
 a^{D_{h-1}}-1\equiv-BD_{h-1}\pmod{B^{h+1}}.
\]
In \eqref{eq:Drec}, the other factor satisfies
$a^{T_{h-1}}\equiv-1\pmod B$, since $T_{h-1}$ is odd. The factor $BD_{h-1}$ already contains $B^h$, so replacing that multiplier by $-1$ is valid modulo $B^{h+1}$. Hence
\[
 D_h\equiv BD_{h-1}\equiv-2B^h\pmod{B^{h+1}}.
\]
This completes the induction.

Because $B\ge4$, the residue $-2$ is nonzero modulo $B$. In a telescoping sum $T_k-T_h$, every term after $D_h$ is divisible by $B^{h+1}$, proving \eqref{eq:sharp-telescope}. The same congruence proves both permanence from height $m$ and failure at every earlier height.
\end{proof}

A useful independent check comes from the local laws. For every odd $q\mid B$,
\[
 v_q(D_h)=h\,v_q(B).
\]
Writing $t=v_2(B)$ and $s=v_2(B-2)$ gives
\[
 v_2(D_h)=s+h(s+t-1).
\]
If $B$ is a power of $2$, then $t\ge2$ and $s=1$, so $v_2(D_h)=1+ht$; even in this case the number of full base-$B$ factors is exactly $h$. These valuation checks recover the exact number of stable digits, while the binomial proof additionally determines the first unstable digit.

\section{Resolving the three OEIS A324017 conjectures}\label{sec:OEIS}

With $B=2n$ and $a=B-1$, the array in OEIS A324017 \cite{OEIS324017} is
\begin{equation}\label{eq:array}
 A(m,n)=\res{T_m(2n-1)}{(2n)^m}.
\end{equation}
The row coordinate is the tower height $m$; the column coordinate is $n$. We work with $n>1$, so $B\ge4$. The exceptional column $n=1$ has base $1$ and is identically $1$.

\begin{lemma}[An exponent period that is sufficient here]\label{lem:period}
For even $B\ge4$, $a=B-1$, and $m\ge1$,
\begin{equation}\label{eq:period}
 a^{B^m}\equiv1\pmod{B^{m+1}}.
\end{equation}
\end{lemma}
\begin{proof}
For an odd prime $q\mid B$, the exponent $B^m$ is even, and odd-prime lifting applied to $a^2$ gives
\[
 v_q(a^{B^m}-1)=v_q(a+1)+v_q(B^m)=(m+1)v_q(B).
\]
For $q=2$, put $t=v_2(B)$ and $s=v_2(B-2)$. Binary lifting gives
\[
 v_2(a^{B^m}-1)=mt+s+t-1.
\]
If $t\ge2$, then $s=1$ and this is $(m+1)t$. If $t=1$, then $s\ge2$, and it is at least $(m+1)t$. Thus every prime-power factor of $B^{m+1}$ divides the expression.
\end{proof}

\subsection{The first conjecture is true}

The first comment asserts, in our notation, that whenever $m,n,k>1$ and $T_m(a)\equiv k\pmod{B^m}$, one has $a^k\equiv k\pmod{B^m}$.

\begin{proposition}\label{prop:oeis1}
For every even $B\ge4$ and $m\ge1$, every nonnegative $k\equiv T_m(B-1)\pmod{B^m}$ satisfies
\[
 (B-1)^k\equiv k\pmod{B^m}.
\]
\end{proposition}
\begin{proof}
Lemma~\ref{lem:period} implies that exponent reduction modulo $B^m$ is valid modulo $B^m$. Consequently
\[
 a^k\equiv a^{T_m}=T_{m+1}\equiv T_m\equiv k\pmod{B^m},
\]
where Theorem~\ref{thm:sharp} gives the middle congruence.
\end{proof}

\subsection{The second conjecture is true}

The second comment concerns the penultimate base-$B$ digit.

\begin{proposition}\label{prop:oeis2}
For $B\ge4$ even and $m\ge2$,
\[
 \floor{\frac{T_m(B-1)}{B}}\equiv B-2\pmod B.
\]
\end{proposition}
\begin{proof}
By Theorem~\ref{thm:sharp}, $T_m\equiv T_2\pmod{B^2}$. Since $B-1$ is odd, the binomial theorem gives
\[
 T_2=(B-1)^{B-1}=-(1-B)^{B-1}
 \equiv -1+B(B-1)\equiv B^2-B-1\pmod{B^2}.
\]
This residue equals $(B-2)B+(B-1)$ and lies between $0$ and $B^2-1$. Its penultimate digit is therefore $B-2$.
\end{proof}

\subsection{The third conjecture is false in every nondegenerate case}

The third comment would identify the height-$m$ residue modulo $B^{m+1}$ with $A(m+1,n)$. Theorem~\ref{thm:sharp} gives the exact opposite conclusion:
\begin{equation}\label{eq:oeisrepair}
 \boxed{\ A(m+1,n)=
 \res{\res{T_m(B-1)}{B^{m+1}}-2B^m}{B^{m+1}}.\ }
\end{equation}
The correction term is nonzero modulo $B^{m+1}$ for every $n>1$. Thus the third assertion fails for every such $n$ and every $m>1$, not merely for an isolated example.

At the smallest allowed pair $m=n=2$,
\[
 T_2(3)=27,\qquad T_3(3)=3^{27}=7625597484987\equiv59\pmod{64}.
\]
The residues are distinct, with difference $59-27=32\equiv-2\cdot4^2\pmod{4^3}$.

For the base $11$, the correct computations are
\[
 T_3(11)\equiv19571\pmod{12^4},\qquad
 T_4(11)\equiv16115\pmod{12^4}.
\]
Their difference is $-3456=-2\cdot12^3$. Also $A(3,6)=563$ and $A(4,6)=16115$. These values help separate the height and column coordinates when checking the illustrative comment in the source entry.

\section{Fixed-point dynamics and a factorization-free algorithm}\label{sec:algorithm}

\begin{theorem}[The finite dynamical system]\label{thm:dynamics}
Let $B\ge4$ be even, $a=B-1$, and $M=B^m$ with $m\ge1$. The map
\[
 f:\Z/M\Z\longrightarrow\Z/M\Z,\qquad f([x])=[a^x]
\]
is well defined and has exactly one fixed point, namely $R_m=A(m,B/2)$. Every odd starting residue reaches $R_m$ after at most $m$ iterations, and every starting residue reaches it after at most $m+1$ iterations. These uniform bounds are sharp, attained by starts $1$ and $0$, respectively. There are no other periodic cycles.
\end{theorem}
\begin{proof}
Lemma~\ref{lem:period} makes $a^x\pmod M$ independent of the chosen representative $x$. Proposition~\ref{prop:oeis1} supplies the fixed point $R_m$, which is odd.

For odd $x$ and $y$, both $a^x$ and $a^y$ are $-1$ modulo $B$. If, more generally, $B^r\mid x-y$ with $r\ge1$, then the lifting identities show that
\begin{equation}\label{eq:contraction}
 B^{r+1}\mid a^x-a^y.
\end{equation}
To check this, interchange $x,y$ if necessary and factor the difference. For an odd $q\mid B$, the valuation gain is $v_q(a+1)=v_q(B)$ because $x-y$ is even. For $q=2$, the gain is $s+t-1\ge t=v_2(B)$, with $s=v_2(B-2)$ and $t=v_2(B)$, exactly as in Lemma~\ref{lem:period}.

Compare the orbit of an odd start with the fixed orbit $R_m$. One iteration gives agreement modulo $B$, and \eqref{eq:contraction} gains another factor of $B$ each time. Agreement modulo $B^m$ follows after $m$ iterations. Any start becomes odd after one iteration, giving $m+1$ in general. A fixed point must itself be odd, so the same argument proves uniqueness. Every periodic orbit must eventually reach $R_m$, hence no other cycle exists.

The orbit from $1$ is $T_0,T_1,\ldots$ and reaches the limiting residue for the first time at height $m$ by Theorem~\ref{thm:sharp}. The orbit from $0$ begins $0,1,T_1,\ldots$, so its first arrival is one step later.
\end{proof}

\subsection{Digit lifting at increasing precision}

There is no need to iterate at the final modulus throughout. Start with $x_1=B-1$ and set
\begin{equation}\label{eq:lift}
 x_{j+1}=\res{(B-1)^{x_j}}{B^{j+1}}\qquad(1\le j<m).
\end{equation}
Then $x_j=R_j$ for every $j$. Indeed, if $x_j\equiv T_j\pmod{B^j}$, the stronger modulus in Lemma~\ref{lem:period} permits replacing the exponent $T_j$ by $x_j$ modulo $B^{j+1}$, giving $x_{j+1}\equiv T_{j+1}\pmod{B^{j+1}}$.

The core implementation is only a few lines:
\begin{lstlisting}[language=Python]
def stable_radix_residue(B: int, digits: int) -> int:
    if B < 4 or B % 2 or digits < 1:
        raise ValueError("B must be even >= 4; digits must be >= 1")
    a, x, modulus = B - 1, B - 1, B
    for _ in range(1, digits):
        modulus *= B
        x = pow(a, x, modulus)
    return x
\end{lstlisting}
The distributed implementation additionally checks integer types and rejects booleans. It uses no factorization and never constructs a power tower. If $L=\lceil m\log_2B\rceil$ is the output bit length and $\mathsf M(L)$ is the cost of $L$-bit multiplication, ordinary binary modular exponentiation gives the coarse upper bound
\[
 O\!\left(\sum_{j=1}^m (j\log B)\mathsf M(j\log B)\right)
 \subseteq O(mL\mathsf M(L)).
\]
The algorithm stores $O(L)$ bits, apart from arithmetic workspace and optional output tables. This is a bound on the displayed algorithm, not an optimality claim.

\subsection{Why naive exponent reduction needs a guard}

For a generic pair $(a,M)$, iterating $r\leftarrow a^r\bmod M$ need not compute true tower residues. For example, with $a=2$, $M=7$, this iteration from $r_0=1$ gives $1,2,4,2,4,\ldots$, while $T_4(2)=65536\equiv2\pmod7$, not $4$.

The included function \texttt{residue\_towers} therefore checks the sufficient condition $a^M\equiv1\pmod M$ before using residue iteration. This condition holds for $a=B-1$, $M=B^m$ by Lemma~\ref{lem:period}. It also holds for every $a\equiv9\pmod{10}$ and $M=10^k$: binary lifting provides the required factors of $2$, and odd-prime lifting applied to $a^2$ provides the factors of $5$.

A separate reference routine uses an Euler-totient chain for moduli supported on $2$ and $5$. It reduces exponents modulo $\varphi(M)$, not modulo $M$. Since bases coprime to $10$ remain units at every modulus in that chain, ordinary Euler reduction is valid throughout. The two implementations are cross-checked against each other.

\section{A consequence for higher up-arrow operations}\label{sec:arrows}

The arithmetic result also survives replacement of tower height by faster-growing operations. To avoid ambiguity about conventions, define
\[
 U_1(a,n)=a^n,\qquad U_r(a,0)=1,
\]
and, for $r\ge2$,
\begin{equation}\label{eq:up-arrow-def}
 U_r(a,n+1)=U_{r-1}(a,U_r(a,n)).
\end{equation}
Thus $U_2(a,n)=T_n(a)$; $U_r$ is the $r$-arrow operation under this convention.

\begin{theorem}[Exact reduction of higher-arrow thresholds]\label{thm:arrows}
Let $B\ge4$ be even, $a=B-1$, and $R_m=A(m,B/2)$. For every $r\ge2$, there is a least $\sigma_r(m)\ge1$ such that
\[
 U_r(a,n)\equiv R_m\pmod{B^m}
 \quad\Longleftrightarrow\quad n\ge\sigma_r(m).
\]
These thresholds satisfy
\begin{align}
 \sigma_2(m)&=m,\label{eq:sigma2}\\
 \sigma_r(m)&=1+\min\{j\ge0:\ U_r(a,j)\ge\sigma_{r-1}(m)\}
 \quad(r\ge3),\label{eq:sigmarec}\\
 \sigma_r(m)&\le m.\label{eq:sigmabound}
\end{align}
In particular, for every $r\ge2$ and $n\ge m$, the residue is the same $R_m$, independently of the arrow rank.
\end{theorem}
\begin{proof}
Theorem~\ref{thm:sharp} gives the exact equivalence and threshold when $r=2$. The integer-valued functions in \eqref{eq:up-arrow-def} are strictly increasing in $n$ and satisfy $U_r(a,n)\ge n+1$. These facts follow inductively from $a^n\ge n+1$ for $a\ge3$ and the defining recurrence.

Assume the assertion for rank $r-1$. For $n\ge1$,
\[
 U_r(a,n)=U_{r-1}(a,U_r(a,n-1))
\]
and the right-hand side has the desired residue exactly when
$U_r(a,n-1)\ge\sigma_{r-1}(m)$. The set of such $n$ is a terminal interval because $U_r(a,\cdot)$ is strictly increasing. At $n=0$, the value $1$ is not congruent to $R_m$, since $R_m\equiv-1\pmod B$ and $B\ge4$. This proves \eqref{eq:sigmarec} and the exact equivalence.

Finally, by induction $\sigma_{r-1}(m)\le m$, while $U_r(a,m-1)\ge m$. Thus the minimum in \eqref{eq:sigmarec} is at most $m-1$, proving \eqref{eq:sigmabound}.
\end{proof}

This is a finite-modulus result for a specified family of bases. It does not construct real or complex tetration, settle uniqueness questions for fractional iteration, or analyze arbitrary long Conway chains or ordinal-indexed fast-growing hierarchies. Its significance here is narrower: a common stable modular residue persists across all finite arrow ranks, even though the ordinary integer sizes change radically.

\section{Reproducibility, proof audit, and remaining questions}\label{sec:repro}

\subsection{The accompanying artifacts}

The archive includes this \LaTeX{} source and its compiled PDF, a standard-library Python module, a verification and data-generation script, exact CSV tables, a JSON primality certificate, and a machine-readable verification report. The principal commands, run from the archive directory, are
\begin{lstlisting}[language=bash]
python code/verify.py
pdflatex -interaction=nonstopmode -halt-on-error article.tex
pdflatex -interaction=nonstopmode -halt-on-error article.tex
\end{lstlisting}
Python 3.10 or later is sufficient for the code; the recorded run used Python 3.13.5. The test run completed with \textbf{118,271 exact assertions passed}. No floating-point arithmetic is used in the verification program.

The main tests cover 2,000 bases ending in 9, at heights $0$ through $10$; 5,150 cross-checks against the independent Euler-chain evaluator; the sharp radix congruence for every even $B$ from $4$ through $100$; exact digit lifting; exhaustive finite dynamics at small moduli; the general-radix local laws; primality and minimality certificates; and the exact valuations in the CRT construction. A sieve finds exactly 23 prime counterexamples below $100000$, beginning
\[
 2749,\ 5749,\ 9749,\ 15749,\ 16249,\ 17749,\ 20749,\ldots.
\]
These finite results are useful regression checks, not evidence sufficient on their own for the infinitude or density theorems.

\subsection{Potential failure points and how they are addressed}

\paragraph{Height conventions.}
The comparison with the source uses $V(a,h)=C_a(h)-C_a(h-1)$ explicitly. At the counterexample, it is $V(2749,2)$ that differs from the eventual speed.

\paragraph{A decimal last digit does not determine the binary class.}
A number ending in 9 need not be $3$ modulo $4$. In fact all failures occur in the other class. Replacing $s$ by $1$ without checking $a\bmod4$ would erase precisely the examples needed to refute the conjecture.

\paragraph{A pairwise congruence need not automatically be permanent.}
Permanence is proved by strict growth of each local valuation and the unique-minimum rule in a telescoping sum. No numerical agreement is promoted to permanence without this proof.

\paragraph{Modular residue zeros can hide insufficient precision.}
The checker chooses decimal precision strictly larger than both predicted valuations. It rejects zero differences before calling its finite-valued valuation routine. The 16-digit columns in the CSV trace are for compact display only; their exact valuation columns were checked at 100-digit precision.

\paragraph{Prime infinitude is not supplied by CRT alone.}
The residue class is proved coprime to its modulus before applying Dirichlet's theorem. The proof of infinitely many primes is a mathematical dependency, not a computational test.

\paragraph{Natural density requires control of an infinite union.}
The proof includes a vanishing-density tail bound. The reported density is for integers in a progression; it is not inferred from a finite prime sample.

\paragraph{Novelty and formal verification.}
The manuscript is not peer reviewed and is not formalized in a proof assistant. Its elementary proofs are intended to be directly checkable. Neither the source audit nor the Python suite establishes priority, exhaustiveness of the literature search, or mechanized proof correctness.

\subsection{Concrete directions left open by this note}

The exact formula eliminates the selected height-2 conjecture and completely classifies the ending-9 case. Further research can now ask sharper questions rather than repeat that conjecture.

\begin{question}[Least prime at each prescribed onset]
Let $P_H$ be the least prime ending in 9 with $H(P_H)=H$, for $H\ge3$. Determine useful upper and lower bounds for $P_H$, and identify which admissible pairs $(s,e)$ achieve the minimum. The theorem gives infinitude for every $H$, but the displayed progression need not contain its smallest example.
\end{question}

\begin{question}[Prime distribution of delays]
Determine an appropriately justified analogue of \eqref{eq:density} when sampling only primes ending in 9. This requires prime-distribution input beyond the infinitude theorem used here and careful treatment of the unbounded valuation conditions.
\end{question}

\begin{question}[Prime factors outside the supported radix]
For general coprime pairs $(a,R)$, describe exact onset laws when some prime divisor of $R$ does not divide $a^2-1$. Multiplicative orders can then introduce additional initial phases. Theorem~\ref{thm:general-radix} deliberately does not cover those primes.
\end{question}

These are directions not resolved by this manuscript; their novelty or open status in the full literature is not asserted. The completed result is the refutation and exact replacement of the selected conjecture, with the accompanying prime-family, density, and radix theorems.

\appendix
\section{A compact certificate of the main counterexample}\label{app:certificate}

For a reader interested only in checking the disproof, the following chain contains the whole mathematical certificate.

The integer $2749$ is prime by trial division through $47$. Define $T_0=1$, $T_{h+1}=2749^{T_h}$, and $D_h=T_{h+1}-T_h$. Since
\[
 v_2(2748)=2,\qquad v_2(2750)=1,\qquad v_5(2750)=3,
\]
we have $v_2(D_0)=2$ and $v_5(D_0)=0$. Every $D_h$ is even, and
\[
 D_h=2749^{T_{h-1}}(2749^{D_{h-1}}-1)\quad(h\ge1).
\]
Binary and odd-prime lifting therefore give
\[
 v_2(D_h)=v_2(D_{h-1})+2,\qquad
 v_5(D_h)=v_5(D_{h-1})+3.
\]
Thus
\[
 v_2(D_h)=2+2h,\qquad v_5(D_h)=3h,
 \qquad C(h)=\min(2+2h,3h).
\]
The speed at height 2 is $C(2)-C(1)=6-3=3$, whereas at every height $h\ge3$ it is $2$. Hence the assertion that all prime bases ending in 9 have permanent speed from height 2 is false.

The alternative Pocklington data in the JSON file give a shorter machine-checkable primality certificate: $229$ is prime, $229\mid2748$, $229>\sqrt{2749}$,
\[
 2^{2748}\equiv1\pmod{2749},\qquad
 2^{12}\equiv1347\pmod{2749},\qquad \gcd(1346,2749)=1.
\]
This optional certificate is not needed for the article's proof, which already uses elementary trial division.

\clearpage
\section{Notation and dependency summary}

\begin{longtable}{p{0.20\textwidth}p{0.73\textwidth}}
\toprule
Symbol&Meaning\\
\midrule
\endhead
$T_h(a)$&Height-$h$ tower, with $T_0(a)=1$.\\
$D_h(a)$&$T_{h+1}(a)-T_h(a)$.\\
$v_q(N)$&Prime-adic valuation of a nonzero integer.\\
$C_a(h)$&Exact stable decimal digit count for a base ending in 9.\\
$V(a,h)$&Newly gained digits: $C_a(h)-C_a(h-1)$.\\
$V(a)$&Eventual value of the speed.\\
$H(a)$&First height from which the speed is always $V(a)$.\\
$s,t,u,e$&$v_2(a-1),v_2(a+1),s+t-1,v_5(a+1)$, respectively.\\
$C_R(h)$&Exact stable count in a supported radix $R$.\\
$J_R(m)$&First permanent height modulo $R^m$; height 0 is allowed.\\
$B$&Even radix, at least 4, with tower base $B-1$.\\
$A(m,n)$&OEIS array $[T_m(2n-1)]_{(2n)^m}$.\\
$R_m$&Unique fixed residue for base $B-1$ modulo $B^m$.\\
$U_r(a,n)$&The explicitly defined $r$-arrow operation.\\
$\sigma_r(m)$&Exact onset of the common residue for $U_r(a,n)$ modulo $B^m$.\\
\bottomrule
\end{longtable}

The logical dependency is short. The two elementary lifting lemmas imply the local affine laws. Those laws imply permanence, the exact decimal transient, minimality of $2749$, and the general-radix formulas. The CRT argument plus Dirichlet supplies the infinite prime families. The density proof adds only elementary residue-class counting and a tail estimate. A separate binomial induction proves the sharper base-$(B-1)$ congruence; exponent-period lifting then supplies the OEIS fixed-point identity and the digit-lifting algorithm. Higher-arrow thresholds follow from the exact tetration threshold and the defining recursion.

\begingroup
\small
\setlength{\parskip}{0pt}
\begin{thebibliography}{9}
\setlength{\itemsep}{3pt}

\bibitem{Ripa2021}
Marco Rip\`a,
\emph{The congruence speed formula},
Notes on Number Theory and Discrete Mathematics \textbf{27} (2021), no.~4, 43--61.
DOI: \href{https://doi.org/10.7546/nntdm.2021.27.4.43-61}{10.7546/nntdm.2021.27.4.43-61}.
Journal version: \url{https://nntdm.net/papers/nntdm-27/NNTDM-27-4-043-061.pdf}.
Author version: \href{https://arxiv.org/abs/2208.02622v1}{arXiv:2208.02622v1} (2022).
The selected claim is Conjecture 2, journal p.~57, and Conjecture 4.1, arXiv p.~14. Versions examined September 19, 2026.

\bibitem{RipaOnnis2022}
Marco Rip\`a and Luca Onnis,
\emph{Number of stable digits of any integer tetration},
Notes on Number Theory and Discrete Mathematics \textbf{28} (2022), no.~3, 441--457.
DOI: \href{https://doi.org/10.7546/nntdm.2022.28.3.441-457}{10.7546/nntdm.2022.28.3.441-457}.
Author version: \href{https://arxiv.org/abs/2210.07956v1}{arXiv:2210.07956v1}.
See especially Corollary 2.2 and Section 2.2 for prior eventual-speed formulas and transient bounds.
Examined September 19, 2026.

\bibitem{OEIS324017}
OEIS Foundation Inc.,
\emph{The On-Line Encyclopedia of Integer Sequences}, entry A324017,
created by Davis Smith, March 28, 2019.
\url{https://oeis.org/A324017}.
Definition, conjectural comments, and coordinate convention examined September 19, 2026.

\bibitem{Dirichlet}
P.~G.~Lejeune Dirichlet,
\emph{There are infinitely many prime numbers in all arithmetic progressions with first term and difference coprime},
original paper published in 1837; English translation by Ralf Stephan,
\href{https://arxiv.org/abs/0808.1408v2}{arXiv:0808.1408v2} (2014).
Used only for infinitude of primes in a reduced arithmetic progression.

\end{thebibliography}
\endgroup
\end{document}
